\minitopic{“有价值”至少包含四个不同问题}讨论知识的价值时，首先要固定比较对象。第一，真信念是否比假信念更有工具价值，例如帮助我们到达目的地？第二，知识是否比同内容的偶然真信念更好？第三，获得知识的过程是否表现了主体能力，因而值得把成功归功于主体？第四，一种认识成就是否还能支持新的探究、解释和共同信任？如果不区分这些问题，“知识当然有用”和“知识具有独特价值”会被误当作同一结论。

淹没问题集中攻击第二问：若真值已经带来全部结果价值，可靠过程似乎像一杯好咖啡由一台可靠机器冲出一样，不会让成品更好。可靠主义可以回答，可靠来源使真信念更稳定、更容易在相似场景中保持，也让他人有理由继续信任该来源；德性认识论则强调成功因能力而成，具有成就价值\parencite{zagzebski1996,sosa2007}。但二者不能只换一个赞美词。稳定性需要说明跨哪些变化，成就需要说明主体贡献为何足够显著；否则一个完全依赖他人的听者或自动系统会使功劳归属再次成为问题。

\minitopic{次级价值与更高阶价值}知识不仅可以作为一个已经完成的状态被评价，还可以从“值得追求”和“形成更大结构”两个层面评价。可靠探究值得选择，即使一次偶然猜测恰好更快得到真相；一组彼此连接的知识又可能形成理解，使主体能够回答为什么、预测变化并发现异常。后一种价值不能简单还原为更多真命题，因为同样数量的孤立事实未必让人把握依赖关系\parencite{kvanvig2003,grimm2006}。

这也解释为什么教育不能只按正确答案评分。若学生知道标准答案，却不知道哪些证据支持它、条件改变后结论是否保留，他拥有的认识成果十分脆弱。反过来，一个学生使用有意简化的模型，却能指出理想化位置、推出反事实后果并在预测失败时修订，可能具有更深理解。评价不应因此放弃真理约束，而应检查局部失真是否服务于更准确的关系把握。

\begin{center}
\begin{tabularx}{\textwidth}{p{0.16\textwidth}XXX}
\toprule
表现 & 知道事实 & 理解关系 & 表现智慧\\
\midrule
背诵者 & 可复述结论 & 不能迁移或解释 & 不知何时暂停判断\\
模型使用者 & 掌握主要事实 & 能作反事实推演 & 知道模型边界并寻求复核\\
成熟探究者 & 能追踪证据 & 能整合多种解释 & 能选择值得追问的问题并修订目标\\
\bottomrule
\end{tabularx}
\end{center}

\minitopic{理解是否必须完全无误}有些科学模型包含理想化假设，例如无摩擦平面或完全理性行动者，却仍能带来理解。这似乎支持“理解可以建立在假命题上”。但从局部理想化直接推出理解无需真理约束过快。应区分三件事：模型字面描述是否完全真实，模型保留的依赖关系是否大体正确，主体是否知道理想化会在哪里失效。一个模型可以在第一层为假，在后二层仍可靠；若主体把理想化当作现实全部结构，理解就会崩溃。

因此，理解的检验不只是要求主体讲出一个流畅故事。至少还要看它能否支持新情境迁移、反事实推演、解释比较和错误定位；竞争解释产生不同预测时，主体能否说明什么证据会区分它们。叙事连贯本身很廉价，阴谋论也能连贯。理解需要让世界对解释施加摩擦。

\minitopic{个人成就与共同认识善可能冲突}以个人功劳评价认识成就，容易低估证言、实验团队、数据库和公共机构的贡献；只评价系统总准确率，又可能容忍少数人的发言被系统性忽略。共同认识善至少包含真信息的生产、错误的公开纠正、解释资源的可获得、异议者的参与机会以及信用的合理分配。它们可能冲突：快速形成共识有利于行动，却可能压低异议；开放全部材料有利于复核，却可能损害隐私。

智慧并不是把这些价值加成一个总分，而是在具体角色中决定哪些问题值得优先、哪些损失不能被平均收益抵消，以及何时需要修改目标。个人可以拥有大量知识而在问题选择上十分愚蠢；机构也可以平均准确却反复把错误成本转嫁给同一群体。认识价值理论若不能进入教育、专家制度和公共决策，就没有完成从“好答案”到“好的认识生活”的解释。

\begin{argumentbox}
评价一项认识成果时依次追问：它是否为真；是否来自能抵抗相关错误的方法；成功能否合理归功于主体或团队；能否支持解释、迁移和继续探究；成果及其成本如何在人群中分配。任何单一层面的高分都不能自动替代其余层面。
\end{argumentbox}

\index{次级认识价值}
\index{理解的真理约束}
\index{共同认识善}
